% Created 2018-10-18 Thu 11:23
% Intended LaTeX compiler: pdflatex
\documentclass[11pt]{article}
\usepackage[utf8]{inputenc}
\usepackage[T1]{fontenc}
\usepackage{graphicx}
\usepackage{grffile}
\usepackage{longtable}
\usepackage{wrapfig}
\usepackage{rotating}
\usepackage[normalem]{ulem}
\usepackage{amsmath}
\usepackage{textcomp}
\usepackage{amssymb}
\usepackage{capt-of}
\usepackage{hyperref}
\pagenumbering{roman}
\usepackage[margin=0.80in]{geometry}
\usepackage[backend=biber, style=ieee, url=false]{biblatex}
\usepackage{float}
\usepackage[super,negative]{nth}
\usepackage[capitalise]{cleveref}
\usepackage{pst-node,transparent,ragged2e}
\addbibresource{/home/nlisa/.spacemacs.d/org-files/bibliography.bib}
\addbibresource{bibliography.bib}
\DeclareFieldFormat[inproceedings]{citetitle}{\textit{#1}}
\DeclareFieldFormat[inproceedings]{title}{\textit{#1}}
\DeclareFieldFormat[misc]{citetitle}{#1}
\DeclareFieldFormat[misc]{title}{#1}
\renewcommand*{\bibpagespunct}{%
\ifentrytype{inproceedings}
{\addspace}
{\addcomma\space}}
\AtEveryCitekey{\ifuseauthor{}{\clearname{author}}}
\AtEveryBibitem{\ifuseauthor{}{\clearname{author}}}
\SpecialCoor
\author{Nyameko Lisa}
\date{30 September 2018}
\title{IPSP078 - Legal Aspects of Electronic Commerce\\\medskip
\large Assignment 2: 724205}
\hypersetup{
 pdfauthor={Nyameko Lisa},
 pdftitle={IPSP078 - Legal Aspects of Electronic Commerce},
 pdfkeywords={},
 pdfsubject={},
 pdfcreator={Emacs 26.1 (Org mode 9.1.14)},
 pdflang={English}}
\begin{document}

\maketitle
\addvspace{110pt}
\centering{
\pnode(0.5\textwidth,-0.5\textheight){thisCenter}
\rput(thisCenter){%\transparent{0.25}
\includegraphics[width=2.7in]{/home/nlisa/course/llb/wipo-unisa/UNISACoatofArms.eps}}}

\justifying
\addvspace{110pt}
\section*{Declaration}
\label{sec:orgd407c46}
I know that plagiarism is to use someone else’s work and pass it off as my own.\\
I know that plagiarism is wrong.\\
I confirm that this assignment is my own work.\\
I have acknowledged in the bibliography accompanying the assignment all the sources that I have used.\\
I have not directly copied without acknowledgement anything from the Internet or from any other source.\\
I have indicated every quotation and citation in a footnote or bracket linked to that quotation.\\
I have not allowed anyone else to copy my work and to pass it off as their own work.\\
I understand that if any unacknowledged copying whatsoever appears in my assignment I will receive zero per cent for the assignment.\\
I am aware of the UNISA policy on plagiarism and understand that disciplinary proceedings can be instituted against me by UNISA if I contravene this policy.\\
I indicate my understanding and acceptance of this declaration by
entering my name here under:
\begin{itemize}
\item Name: \textbf{Nyameko Lisa} (Student Number: \textbf{7874-909-3})
\end{itemize}

\subsection*{NOTE}
\label{sec:orgb8c8652}
Please note that footnotes will be denoted as \footnote{This is a footnote.} and will
appear at the bottom of the page.\\
References will be denoted by \cite{Cachalia15_spring_forest_v_wilberry} and will appear at the end of the document.
\newpage

\pagenumbering{arabic}

\section{Question 1 [25]}
\label{sec:orga74728f}

The following discussions pertain to \citetitle{Cachalia15_spring_forest_v_wilberry}
\cite{Cachalia15_spring_forest_v_wilberry},

\subsection{Briefly state the facts of the case. (5)}
\label{sec:org785f689}

This case is an appeal from a decision of the high court, KwaZulu-Natal, Durban
(Judge Madondo) granting interim \emph{`pendente lite relief'} following an urgent
application by Wilberry (Pty) Ltd t/a as Ecowash \textbf{the respondent}, against Spring
Forest Trading 599 CC \textbf{the appellant}, and Combined Motor Holdings Limited t/a
The Green Machine (CMH)\footnote{Cited here as it has an interest in the relief
claimed, but is not party to the present appeal.}.\\

The appeal concerns a series of emails purporting consensual written
cancellation of several agreements in which the applicant leased
MDU's\footnote{Mobile Dispensing Units.} from the respondent, which the applicant
required for use in its mobile car wash business. The agreement included a
non-variation clause stipulating that changes to or cancellation of the
agreement would only be effective if recorded in writing and signed by both
parties, and it also afforded the appellant the right to promote, operate and lease out
the respondent's MDU's to third parties.\\

The appellant subsequently reneged on its rental obligations in terms of the
agreement and after negotiations the parties had reached a verbal understanding
to cancel the standing agreements. The terms of the cancellation of the
agreements, were recorded in a series of email correspondences between the
respective representative(s) for the respondent Mr Nigel Keirby-Smith and those
for the appellant, Mr Gregory Stuart Hamilton and Mr Walter Burger. These
included the settlement of the rental payment in arrears owed by the appellant
and the return of the MDU's to the respondent.

\subsection{Identify and discuss the principle of law in the case. (5)}
\label{sec:orgb7b1d14}

The appeal concerned a series of emails purporting to the consensual
cancellation of the written agreements. Stipulated within these written
agreements was that `consensual cancellation' to be effected \textbf{in writing and
signed by both parties}. As per the provisions of [sections 11, 12 and
13]\cite{rsa02_elect_comm_trans_act}, statutory regulations of the Republic
afford legal recognition to transactions concluded electronically via email. The
SCA\footnote{Supreme Court of Appeals.} was required to consider whether the dispute
arising from the exchange of emails between the parties, did indeed satisfy the
\textbf{writing} [section 12]\cite{rsa02_elect_comm_trans_act} and \textbf{signature} [section
13]\cite{rsa02_elect_comm_trans_act} requirements, thereby constituting legal
consensual cancellation.\\

The appellant settled the rental fees which were, but continued its car washing
operations at the locations of the rental agreements, arguing that it was
entitled to do so as the agreements between itself and the respondent had been
effectively canceled. The appellant then entered into subsequent agreements with
another entity to conduct the same operations. To which the respondent denied
the validity of the cancellations to the agreements, and sought interdictory
relief preventing the appellant from conducting said operations, protecting its
proprietary rights in it's MDU's pending the institution of breach of
agreement actions.\\

In the original action the HC\footnote{(Durban) High Court.} found that the email
communication constituted inconclusive negotiations between the parties and did
not evince and intention to cancel the agreements, adding that the parties had
not explicitly specified that their agreements could be canceled via email.\\

The respondents argued that the emails did not comply with the requirements of
[section 13(1)]\cite{rsa02_elect_comm_trans_act}, lacking an endorsement of an
`advanced electronic signature'. Moreover, the respondent also argued that the
appellant cannot bring the emails into the ambit of the law as per the
provisions of [section 13(3)]\cite{rsa02_elect_comm_trans_act}, contending that
the emails pertain to oral negotiations about the agreement and hence cannot
constitute separate electronic transactions, nor were there any controls
implemented to identify the parties and indicate their approval.\\


\subsection{What was the court's decision regarding the issues raised? (5)}
\label{sec:org6a3ce6c}

On appeal the SCA upheld that it was not in dispute whether the emails satisfied
the requirement that the cancellation need be `in writing', as per [section
12(a)]\cite{rsa02_elect_comm_trans_act}, of which statutory requirements are
satisfied if said correspondence is in the form of a `data message'. Instead the
SCA held that the real issue was whether the names of the parties at the foot of
their respective emails, indeed constituted signatures as per the provisions of
[section 13(1) and 13(3)]\cite{rsa02_elect_comm_trans_act}.\\

The SCA held that the stipulations for the requirement of an `advanced
electronic signature' [section 13(1)]\cite{rsa02_elect_comm_trans_act} did not
apply to the circumstances of this case, whereas however the less restrictive
requirements of an 'electronic signature' [section
13(3)]\cite{rsa02_elect_comm_trans_act} do indeed apply. Wherein it need only be
demonstrated that the requirement is indeed met if a method of electronic
signature used to identify a party and indicate their approval of the contained
information [section 13(3)(a)]\cite{rsa02_elect_comm_trans_act}, and that the method
was indeed `appropriately reliable' for the intended purposes of the information
communicated [section 13(3)(b)]\cite{rsa02_elect_comm_trans_act}. In the SCA's analysis of [section 13(3)]\cite{rsa02_elect_comm_trans_act}, Justice
Cachali argued [paragraph 19]\cite{Cachalia15_spring_forest_v_wilberry}:
\begin{quote}
\textit{‘The respondent submits that the phrase: “Where the signature of a person is required by law” (emphasis added) in s 13(1) it should be interpreted not only to include formalities required by statute but must also incorporate instances where parties to an agreement impose their own formalities on a contract, as in this case. And, so the contention goes, because the parties required their signatures for the contracts to be cancelled the requirement could only be satisfied by the use of an advanced electronic signature as contemplated in s 13(1), which did not occur in this case.’}
\end{quote}

The SCA ordered that the appeal be upheld with costs, and that the order of the
HC be set aside and dismissed with costs [paragraph 32]\cite{Cachalia15_spring_forest_v_wilberry}.

\subsection{What is the impact of the case on the law? (5)}
\label{sec:org2e71b75}

The SCA held that courts seek to determine whether the method of signature used
fulfills the function of a signature which is to authenticate the identity of a
signatory and to appropriately validate it's authenticity, rather than insisting
on the form a signature may assume. The judgment held that the approach of
courts has in general been pragmatic and not overly formal, citing a case
[paragraph 148F-G]\cite{elec_signa63_putter_v_provincial_insurance} where
the courts have accepted any mark made by a person attesting to a document
[paragraph 25]\cite{Cachalia15_spring_forest_v_wilberry}:
\begin{quote}
\textit{`In the days before electronic communication, the courts were willing to accept any mark made by a person for the purpose of attesting a document, or identifying it as his act, to be a valid signature. They went even further and accepted a mark made by a magistrate for a witness, whose participation went only as far as symbolically touching the magistrate’s pen.'}
\end{quote}

The SCA held that the typed written names of the respective parties appearing at
the foot of the emails in question, were indeed intended to identify the
parties, amounted to data that was logically associated with data in the body of
the emails, and thus constituting legal electronic signatures, [paragraph 29]\cite{Cachalia15_spring_forest_v_wilberry}:
\begin{quote}
\textit{`There is no dispute regarding the reliability of the emails, the accuracy of the information communicated or the identities of the persons who appended their names to the emails. On the contrary, as I have found earlier, the emails clearly and unambiguously evinced an intention by the parties to cancel their agreements. It ill-behoves the respondent, which responded to clear questions by email itself, to now rely on the non-variation clauses to escape the consequences of its commitments made at the meeting on 25 February 2013 which were later confirmed by email.'}
\end{quote}

Given that the SCA upheld the appeal, finding that the email cancellation of the
agreements in question was indeed valid, raises the important and difficult
question of whether courts of the Republic would hold a similar view in terms of
an agreement canceled via other electronic instant text messaging platforms and
social media, such as SMS\footnote{Short Message Services.}, WhatsApp, FaceBook
status updates, Twitter, etc\ldots{}

\subsection{What are the opinions of other authors regarding the principle discussed in the case? (5)}
\label{sec:orgddd3b9f}

With reference to [section 13]\cite{rsa02_elect_comm_trans_act} Shumani L. Gereda
\cite{gereda16_elec_comms_trans_act}, argued that where contracting parties failed
or neglected to agree on the electronic signature to be used, the presumption is
that the requirement of a signature is met if a method is used to identify the
person (who has purportedly `signed' the electronic communication) and to
indicate said person's approval of the data message. Such intent can be made by
any means from which said person's intent can be inferred.\\

In the newsletter for Sonnenberg Property \& Legal Advice \cite{fritz_sonnenberg},
Fritz Sonnenberg advises how the precedent established by the  findings of the
SCA would be applicable to most contracts (including residential and commercial
leases), where electronic communications are used to effect changes. However
electronic signatures cannot be used for long-term\footnote{Exceeding 20 years.}
lease of land, the signing of a will, agreements for the sale of immovable
property and bills of exchange. Adding that it would be beneficial to include a
clause in an agreement stipulating the nature of electronic communications that
will constitute acceptable means of written communications. Lastly he advises
that one should ensure email communications should be signed if one requires
them to be binding, and one should explicitly declare that they are not binding
should one not want them to be bound.\\

Simone Dickson, a director at the firm Cliffe Dekker Hofmeyer provides similar
advise \cite{simone_dickson}, advising contracting parties to be cognizant of the
SCA's judgment and to carefully consider any electronic communications between
themselves. Moreover she advises that parties include express provisions in the
contract which clearly regulate how electronic communications will apply in
order to avoid disputes and ambiguity.\\

Ryan Tucker of RM Tucker Attorneys, also provides similar advise \cite{rmtucker},
clarifying and reiterating that contract negotiators making amendments to, or
canceling agreements via email, or other electronic means, with a simple
signature at the foot of the message will be sufficient to effect such
amendments/cancellations, unless one explicitly imposes more stringent
requirements in their agreements.\\

In additional to the above considerations, attorneys from Abrahams \& Gross
explicitly enumerate \cite{abgross} the laws/transactions to which [section 12 and
13]\cite{rsa02_elect_comm_trans_act} are not applicable, where variations and/or
cancellations must be effected in the traditional manner written on paper and
signed by hand:
\begin{itemize}
\item Wills Act 7 of 1953 (Wills);
\item Alienation of Land Act 68 of of 1981 (An agreement for alienation/sale of immovable property and an agreement for the long terms lease of immovable property in excess of 20 years as provided for in the aforesaid act);
\item Bills of Exchange Act 34 of 1964 (bills of exchange); and
\item Stamp Duties Act 77 of 1968 (stamp duties).
\end{itemize}

\printbibliography
\end{document}